% Stage 8 section framing — 2.4 Radiation Conditions and Uniqueness
\section{Radiation Conditions and Uniqueness}
\label{sec:ch24}
Green-function inversion in Section~\ref{sec:ch23} reconstructs fields from sources on compact support. Section~\ref{sec:ch24} imposes outgoing radiation conditions that separate measurable far-zone content from reactive storage \cite{sommerfeld1949}. A dense near-field map without outgoing filtering can fit non-radiating currents \cite{devaney2004}.


Maxwell's equations alone do not distinguish fields radiating away from a source from
fields arriving from infinity, yet every laboratory measurement assumes the former.
Antennas are fed by finite sources, compact ranges suppress external illumination, and
receivers attribute detected power to identifiable radiators rather than unexplained
incoming waves \cite{sommerfeld1949}. Section~\ref{sec:ch23} supplied Green
kernels that map impressed currents to fields; on exterior domains those kernels remain
ambiguous until an outgoing branch is fixed at infinity \cite{jackson1999}. Retarded
causality narrows the admissible class but does not fully replace Sommerfeld selection
when harmonic solvers invert $\mathcal{L}_\omega$ at fixed $\omega$
\cite{stratton1941}. The subsections below state outgoing radiation conditions, embed them
as resolvent selection, prove uniqueness on the radiative branch, and record
reconstruction limits when that branch is omitted.
